\documentclass[a4paper, 12pt]{article}
\usepackage[spanish]{babel}
\usepackage[hmargin=2cm,vmargin=2.5cm]{geometry}
\usepackage{enumerate}
\usepackage{makecell}
\usepackage{graphicx}
\usepackage{hyperref}
\usepackage{amsmath}
\usepackage[backend=biber,style=apa, url=true, sortcites]{biblatex}
\usepackage[table]{xcolor}
\usepackage{minted}
\usepackage{graphicx}
\usepackage{fancyhdr}  % Agrega el paquete fancyhdr
\usepackage{subcaption}

\addbibresource{references.bib}
\hypersetup{
	colorlinks,
	citecolor=black,
	filecolor=black,
	linkcolor=black,
	urlcolor=black
}

\setlength{\arrayrulewidth}{0.4mm}

\newcommand{\HRule}{\rule{\linewidth}{0.5mm}}

\begin{document}
    \begin{titlepage}
        \begin{center}
            % logo
            \includegraphics[width=0.5\textwidth]{figures/logoUAH.png}~\\[2cm]
            
            \textsc{\Large \\Sistemas de Control Inteligente}\\[2cm]
            
            \HRule \\[0.4cm]
            {\LARGE \bfseries Práctica 2. \\ Control Borroso  \\[0.4cm]}
            \HRule \\[3cm]
            
            \large\textbf{Jorge Revenga Martín de Vidales}\\
            \large\textbf{Ángel Salgado Aldao}\\
            \large\textbf{}\\ Grado en Ingeniería Informática \\ Universidad de Alcalá
            
            \vfill
            
            {\large \today}
        \end{center}
    \end{titlepage}

    % Configura los encabezados y pies de página
    \pagestyle{fancy}
    \fancyhf{} % Limpia todos los encabezados y pies de página actuales
    % Encabezado
    \fancyhead[RO,LE]{\textit{Sistemas de Control Inteligente}}
    \fancyhead[LO,RE]{\textit{Control Borroso}}  % Encabezado izquierda del documento
    % Pie de página
    \fancyfoot[LO,RE]{\textit{Universidad de Alcalá}}
    \fancyfoot[RO,LE]{\thepage}  % Número de página en la esquina inferior derecha
    \newpage
    
    \thispagestyle{plain}
    \tableofcontents
    \newpage

    \part{Diseño de un control borroso de posición para un robot móvil.}
    
    \section{a)}
	
	\begin{figure}[htp!]
		\centering
		\includegraphics[width=0.25\textwidth]{figures/parte1/}
	\end{figure}
	
	
	\subsection{Código}
	\inputminted[fontsize=\scriptsize, linenos, breaklines=true, xleftmargin=0.75cm, frame=lines]{matlab}{code/parte1/}
    
    \subsection{i}
    
    \subsection{ii}
    
    \subsection{iii}
    
    \subsection{iv}
    
    \subsection{v}
    
    \section{b)}
    
    \section{c)}
    
    \section{d)}

                
	\part{Diseño de control borroso de posición con evitación de obstáculos.}
	
	
	\section{a)}
	
	\section{b)}
	
	\section{c)}
	
	\section{d)}
	
	\section{e)}
	
	
	
 \end{document}
